\documentclass[11pt]{article}

\usepackage[margin=1in]{geometry}
\usepackage{amsmath, amssymb, amsthm}
\usepackage{booktabs}
\usepackage{hyperref}
\usepackage{enumitem}

\newtheorem{proposition}{Proposition}
\newtheorem{definition}{Definition}
\newtheorem{corollary}{Corollary}
\newtheorem{remark}{Remark}

\newcommand{\Var}{\operatorname{Var}}
\newcommand{\Cov}{\operatorname{Cov}}
\newcommand{\E}{\operatorname{E}}

\title{Mathematical Foundations of the Qwen3-VL Taste-Annotation\\Evaluation Suite}
\author{}
\date{}

\begin{document}
\maketitle

\begin{abstract}
This note formalizes the statistics used throughout the five-benchmark
evaluation of Qwen3-VL-30B as an automated taste-annotation labeler
(\texttt{qwen\_evaluation\_suite/}). It covers: (1) pairwise Cohen's
$\kappa$ and the ``agreement envelope'' framing used to compare model--human
agreement against human--human agreement; (2) plurality consensus labeling
and the confusion-matrix decomposition used under structurally missing
(conditional) fields; (3) the cosine-similarity relaxation of exact-match
accuracy used for the free-text \texttt{noun} field, with a worked
decomposition matching the observed numbers; (4) classical test theory
applied to rater reliability, proving that self-consistency (test--retest
agreement) is necessary but not sufficient for validity, with the
suite's own ``self-fed'' bias as the canonical counterexample; and (5) a
significance test formalizing the claim that five independent ablations
all failing to move human-fed agency accuracy off $0\%$ constitutes strong
evidence of a genuine perceptual ceiling rather than a fixable prompting
gap.
\end{abstract}

\tableofcontents

\section{Pairwise Cohen's $\kappa$ and the agreement envelope}
\label{sec:kappa}

\subsection{Setup}

Let a \emph{rater} be either a human annotator or the model, and let two
raters $A,B$ each assign one label from a finite category set
$\mathcal{C} = \{c_1,\dots,c_K\}$ (e.g.\ $\mathcal{C} =
\{\text{eating},\,\text{possibly-eating},\,\text{not-eating}\}$ for the
\texttt{scenario} field) to each of $n$ clips they both labeled. Write
$y^A_i, y^B_i \in \mathcal{C}$ for their labels on clip $i$.

\begin{definition}[Observed and chance agreement]
The \emph{observed agreement} is
\[
  p_o(A,B) \;=\; \frac{1}{n}\sum_{i=1}^n \mathbb{1}\!\left[y^A_i = y^B_i\right].
\]
Let $\pi^A_k = \frac{1}{n}\sum_i \mathbb{1}[y^A_i = c_k]$ and
$\pi^B_k = \frac{1}{n}\sum_i \mathbb{1}[y^B_i = c_k]$ be each rater's
marginal frequency for category $c_k$. The \emph{chance agreement} under
the assumption that $A$ and $B$ label independently with these marginals is
\[
  p_e(A,B) \;=\; \sum_{k=1}^K \pi^A_k\, \pi^B_k .
\]
\end{definition}

\begin{definition}[Cohen's $\kappa$]
\[
  \kappa(A,B) \;=\; \frac{p_o(A,B) - p_e(A,B)}{1 - p_e(A,B)} .
\]
\end{definition}

$\kappa = 1$ means perfect agreement beyond chance; $\kappa = 0$ means the
raters agree exactly as often as two independent labelers with the same
marginals would; $\kappa < 0$ means they agree \emph{less} than chance
would predict (systematic disagreement).

\subsection{Pooled multi-rater estimate}

The suite has 11 human raters $H_1,\dots,H_{11}$ (each labeling a subset of
the 506-clip pool, up to 3 raters per clip) plus the model, $\mathrm{Q}$.
Rather than Fleiss' $\kappa$ (which requires every clip to have the same
number of raters), Benchmark~3 uses a \emph{pooled pairwise} estimator: for
a field $f$, collect every $(y^{H_i}_j, y^{H_k}_j)$ pair for every clip $j$
that both $H_i$ and $H_k$ labeled, over all $\binom{11}{2}$ rater pairs,
into one pooled sample $\mathcal{P}_{\mathrm{hh}}$; separately collect every
$(y^{H_i}_j, y^{\mathrm{Q}}_j)$ pair into $\mathcal{P}_{\mathrm{qh}}$. Then
\[
  \kappa_{\mathrm{hh}} = \kappa\big(\mathcal{P}_{\mathrm{hh}}\big), \qquad
  \kappa_{\mathrm{qh}} = \kappa\big(\mathcal{P}_{\mathrm{qh}}\big),
\]
using the same observed/chance-agreement formulas with $\mathcal{P}$ in
place of a single rater pair's clip set.

\begin{remark}[Why pooling is a reasonable but conservative estimator]
Pairs drawn from the same clip (e.g.\ $(H_i,H_k)$ and $(H_i,H_l)$ on clip
$j$) are not independent observations, since they share $y^{H_i}_j$. This
does not bias $\hat\kappa$ but does mean its naive standard error is an
underestimate; the suite reports point estimates only; a clip-clustered
bootstrap would be the appropriate correction if confidence intervals were
needed.
\end{remark}

\subsection{The agreement-envelope test}
\label{subsec:envelope}

The scientific question Benchmark~3 asks is not ``is Qwen accurate?''
in isolation, but:
\begin{quote}
  \emph{Is $\kappa_{\mathrm{qh}}$ (and $p_{o,\mathrm{qh}}$) inside the range
  humans achieve with each other, $\kappa_{\mathrm{hh}}$?}
\end{quote}
This reframes evaluation around a data-dependent ceiling rather than a
fixed target of $100\%$, because for genuinely ambiguous clips even
$\kappa_{\mathrm{hh}} < 1$ is the best achievable score --- no rater,
human or model, can exceed the label's intrinsic annotator disagreement.
Formally, if we posit a noisy-annotator model where each human rater
observes the true (possibly ill-defined) label $Y_i$ through a
rater-specific noise channel, $\kappa_{\mathrm{hh}}$ estimates the
\emph{signal-to-noise ratio of the labeling task itself}, and a rater
(model or human) with $\kappa_{\mathrm{model},h} \approx \kappa_{\mathrm{hh}}$
is, by this criterion, indistinguishable from another human annotator.

\begin{table}[h]
\centering
\begin{tabular}{lrrrr}
\toprule
Field & $\kappa_{\mathrm{hh}}$ & $\kappa_{\mathrm{qh}}$ & $p_{o,\mathrm{hh}}$ & $p_{o,\mathrm{qh}}$\\
\midrule
scenario  & 0.44 & 0.42 & 64.8\% & 64.3\% \\
verb      & 0.62 & 0.26 & 69.5\% & 40.5\% \\
condition & 0.28 & 0.28 & 71.5\% & 76.4\% \\
flavor    & 0.52 & 0.26 & 64.6\% & 44.2\% \\
noun      & 0.55 & 0.48 & 60.7\% & 57.8\% \\
\bottomrule
\end{tabular}
\caption{Pooled-pairwise agreement, Benchmark~3 (full 506-clip pool).}
\end{table}

By this test, \texttt{scenario}, \texttt{condition}, and \texttt{noun} are
inside the human envelope ($\kappa_{\mathrm{qh}} \approx \kappa_{\mathrm{hh}}$
or, for \texttt{condition}, exactly equal), while \texttt{verb} and
\texttt{flavor} fall well outside it ($\kappa_{\mathrm{qh}}$ roughly
$0.4\times$ the human value in both cases) --- the same conclusion drawn
qualitatively from the rater-agreement heatmap, now with a single scalar
per field.

\section{Plurality consensus and confusion under structural missingness}
\label{sec:consensus}

\subsection{Plurality (majority-vote) consensus labels}

For a clip $j$ labeled by raters $H_{i_1},\dots,H_{i_m}$ ($m\ge 2$), let
$v_j(c) = \big|\{r : y^{H_{i_r}}_j = c\}\big|$ count votes for category $c$.

\begin{definition}[Plurality consensus]
Let $c_{(1)}, c_{(2)}$ be the categories with the largest and
second-largest vote counts, $v_{(1)} \ge v_{(2)}$ (ties broken
arbitrarily among tied runners-up). The consensus label is defined only
when the winner is \emph{strict}:
\[
  \hat{Y}_j \;=\;
  \begin{cases}
    c_{(1)} & \text{if } v_{(1)} > v_{(2)}, \\
    \text{undefined} & \text{if } v_{(1)} = v_{(2)}\ \ (\text{exact tie}).
  \end{cases}
\]
\end{definition}

With $m=3$ raters, a tie is possible only when all three raters disagree
($v_{(1)} = v_{(2)} = v_{(3)} = 1$); a $2$-$1$ split always yields a defined
consensus. With $m = 2$, consensus is defined exactly when the two raters
agree ($v_{(1)}=2$). Benchmark~3's accuracy-vs-consensus metric,
\[
  \mathrm{Acc}_{\mathrm{cons}}(f) \;=\;
  \frac{1}{|\mathcal{J}_f|}\sum_{j \in \mathcal{J}_f}
  \mathbb{1}\!\left[y^{\mathrm{Q}}_j = \hat{Y}_j\right], \qquad
  \mathcal{J}_f = \{\, j : \hat{Y}_j \text{ defined for field } f \,\},
\]
is therefore computed over clips where the humans reached a strict
plurality, which structurally excludes the hardest 3-way-split clips and
should be read as an upper bound on how well Qwen matches human judgment on
the pool as a whole.

\subsection{Difficulty transfer as a stratified accuracy}

Partition $\mathcal{J}_f$ by the winning margin $v_{(1)} - v_{(2)}$ (for
$m=3$: \emph{unanimous} when $v_{(1)}=3$, \emph{majority} when $v_{(1)}=2$;
for $m=2$: \emph{agree} when $v_{(1)}=2$, and treat a raw $1$-$1$
\emph{split} as its own stratum where a prediction counts as correct if it
matches \emph{either} rater). The \emph{difficulty-transfer} statistic is
the plurality accuracy computed separately within each stratum,
\[
  \mathrm{Acc}_{\mathrm{cons}}(f \mid \text{stratum } s)
  \;=\; \frac{1}{|\mathcal{J}_f \cap s|}\sum_{j \in \mathcal{J}_f \cap s}
  \mathbb{1}\!\left[y^{\mathrm{Q}}_j = \hat{Y}_j\right].
\]

\begin{proposition}[Monotone difficulty transfer implies shared ambiguity]
\label{prop:difficulty}
Suppose model errors are conditionally independent of human disagreement
given the (unobserved) intrinsic ambiguity $A_j\in[0,1]$ of clip $j$, and
that $A_j$ jointly increases both the probability a human rater deviates
from the modal label and the probability the model deviates from it. Then
$\mathrm{Acc}_{\mathrm{cons}}(f\mid s)$ is monotonically decreasing as $s$
ranges from \emph{unanimous} to \emph{split}. Observing this monotone
decrease (empirically: $84.6\% \to 42.6\%$ for scenario, unanimous vs.\
majority) is evidence \emph{for} the shared-ambiguity model over the
alternative hypothesis that Qwen's errors are i.i.d.\ and uncorrelated with
clip difficulty, since the latter predicts a flat
$\mathrm{Acc}_{\mathrm{cons}}(f\mid s)$ across strata.
\end{proposition}

\subsection{Structural (MNAR) missingness of conditional fields}

The four fields \texttt{verb}, \texttt{additional\_condition}, \texttt{noun},
\texttt{flavor\_profile} are only elicited from a rater when that rater's
own \texttt{scenario} answer is not \texttt{not-eating}. This is missingness
\emph{not at random} (MNAR) in the classical taxonomy: whether $y^{r}_{j,f}$
is observed for rater $r$, field $f$, clip $j$ is a deterministic function
of that same rater's \texttt{scenario} label, $R^r_{j} =
\mathbb{1}[y^r_{j,\text{scenario}} \ne \text{not-eating}]$, which is exactly
the mechanism that generates the field values themselves.

\begin{definition}[Comparable-pair restriction]
For a conditional field $f$ and a pair of raters $(A,B)$ on clip $j$, the
pair contributes to the pooled comparison only when \emph{both} raters
observed the field:
\[
  (y^A_{j,f}, y^B_{j,f}) \in \mathcal{P}_f
  \iff R^A_j = 1 \ \text{and}\ R^B_j = 1 .
\]
\end{definition}

\begin{remark}[Why this restriction is necessary, and what it costs]
Comparing on $R^A_j = R^B_j = 1$ only is required because a blank value
means ``not elicited,'' not ``elicited and equal to the empty
string''--- treating it as its own category would inflate agreement
whenever two raters agree on \texttt{scenario} being \texttt{not-eating}
(the fields are trivially equal by construction, contributing no
information about the conditional field itself). The cost is a
selection effect: $\mathcal{P}_f$ is implicitly conditioned on
\emph{scenario agreement} between $A$ and $B$, so $\kappa_{\mathrm{qh}}$
for \texttt{verb}/\texttt{flavor}/\texttt{noun}/\texttt{condition}
already excludes the subset of clips where Qwen and the human disagree
about whether eating is happening at all --- the reported $n$ for these
fields ($n{=}432$ for Qwen$\leftrightarrow$human vs.\ $n{=}347$ for
human$\leftrightarrow$human, Benchmark~3) reflects exactly this
conditioning, not a data-loss artifact.
\end{remark}

\section{Cosine similarity as a soft relaxation of exact-match accuracy}
\label{sec:cosine}

\subsection{Embedding cosine similarity}

The \texttt{noun} field is free text (``cereal'', ``sippy cup'',
``baked beans'', \dots), so exact-string accuracy penalizes correct-but-
differently-worded answers as harshly as wrong ones. Let
$e:\mathcal{W}\to\mathbb{R}^d$ be a sentence-embedding map (MiniLM,
$d=384$) with $\lVert e(w)\rVert_2 = 1$ for all $w$ (embeddings are
L2-normalized before comparison). For ground-truth noun $g_i$ and
predicted noun $\hat{g}_i$ on clip $i$,
\[
  \mathrm{sim}(g_i,\hat{g}_i) \;=\; \cos\theta_i \;=\;
  \frac{e(g_i)\cdot e(\hat{g}_i)}{\lVert e(g_i)\rVert\,\lVert e(\hat{g}_i)\rVert}
  \;=\; e(g_i)\cdot e(\hat{g}_i) \in [-1,1],
\]
the last equality using unit normalization. Because $e$ is trained so that
semantically related words are embedded nearby, $\mathrm{sim}$ acts as a
continuous relaxation of the indicator $\mathbb{1}[g_i = \hat{g}_i]$:
identical strings give $\mathrm{sim}=1$ exactly, while near-synonyms
(``sippy cup'' vs.\ ``bottle'') give intermediate values and unrelated
words give values near $0$ or below.

\subsection{Soft accuracy and its relation to exact-match accuracy}

\begin{definition}[Soft accuracy at threshold $\tau$]
\[
  \mathrm{Acc}_{\mathrm{soft}}(\tau) \;=\;
  \frac{1}{n}\sum_{i=1}^n \mathbb{1}\!\left[\mathrm{sim}(g_i,\hat{g}_i) \ge \tau\right].
\]
\end{definition}
Exact-match accuracy is the degenerate case $\mathrm{Acc}_{\mathrm{exact}}
= \mathrm{Acc}_{\mathrm{soft}}(1)$ under an embedding where distinct
strings almost never map to identical vectors, and the mean similarity
$\overline{\mathrm{sim}} = \frac{1}{n}\sum_i \mathrm{sim}(g_i,\hat{g}_i)$
reported throughout the suite is not a threshold accuracy but a
continuous summary of the same quantity.

\begin{proposition}[Decomposition of mean similarity by exact-match status]
\label{prop:decomp}
Let $\mathcal{M} = \{i : g_i = \hat{g}_i\}$ (exact matches, contributing
$\mathrm{sim}_i = 1$) and $\mathcal{M}^c$ its complement, with
$m = |\mathcal{M}^c|/n$ the mean similarity restricted to mismatches. Then
\[
  \overline{\mathrm{sim}}
  = \mathrm{Acc}_{\mathrm{exact}} \cdot 1 \;+\; (1-\mathrm{Acc}_{\mathrm{exact}})\cdot \bar{m},
  \qquad
  \bar{m} \;=\; \frac{\overline{\mathrm{sim}} - \mathrm{Acc}_{\mathrm{exact}}}{1 - \mathrm{Acc}_{\mathrm{exact}}}.
\]
\end{proposition}
\begin{proof}
Split the sum defining $\overline{\mathrm{sim}}$ over $\mathcal{M}$ and
$\mathcal{M}^c$; every term in $\mathcal{M}$ equals $1$ by definition of
exact match, and $|\mathcal{M}|/n = \mathrm{Acc}_{\mathrm{exact}}$.
Rearranging gives $\bar m$.
\end{proof}

\subsection{Worked example (Benchmark 1)}

Benchmark~1 reports, on the annotator-5 eval set restricted to human
eating/possibly-eating clips, $\mathrm{Acc}_{\mathrm{exact}} = 70.8\%$
($n=48$) and $\overline{\mathrm{sim}} = 82.0\%$ ($n=46$; two additional
clips were excluded where the human noun was blank, so the identity below
is approximate rather than exact over an identical index set). Applying
Proposition~\ref{prop:decomp} with these figures,
\[
  \bar m \;\approx\; \frac{0.820 - 0.708}{1 - 0.708} \;\approx\; 0.384,
\]
i.e.\ on the $\approx 29\%$ of clips where Qwen's noun did \emph{not}
exactly match the human's, the two words were on average only weakly
related (cosine $\approx 0.38$) rather than close synonyms. Cross-checking
against the raw pairs reported in the suite (e.g.\ \texttt{cereal
$\to$ (none)} at $0.13$, \texttt{banana $\to$ bottle} at $0.42$,
\texttt{beans $\to$ baked beans} at $0.81$) confirms the mismatches are a
mix of outright misses (Qwen abstaining or picking an unrelated object)
and near-misses (a related but non-identical object), consistent with
$\bar m$ sitting well below $1$ but above $0$.

\begin{remark}[Why $\overline{\mathrm{sim}} \ge \mathrm{Acc}_{\mathrm{exact}}$
is not automatic]
Proposition~\ref{prop:decomp} shows $\overline{\mathrm{sim}} >
\mathrm{Acc}_{\mathrm{exact}}$ exactly when $\bar m > 0$, i.e.\ when the
model's \emph{mistakes} are, on average, still semantically closer to the
truth than two independently drawn nouns would be. This is an empirical
property of the model's error distribution, not a mathematical
guarantee --- a model that fails by naming a semantically unrelated object
every time it is wrong would have $\bar m \le 0$ and could show
$\overline{\mathrm{sim}} < \mathrm{Acc}_{\mathrm{exact}}$.
\end{remark}

\section{Classical test theory: reliability is not validity}
\label{sec:reliability}

\subsection{Reliability as test--retest agreement}

Classical test theory models an observed measurement as
$X = T + E$, a true-score component $T$ plus measurement error $E$
with $\E[E]=0$ and $\Cov(T,E)=0$. The \emph{reliability} of the
measurement procedure is
\[
  \rho_{XX'} \;=\; \frac{\Var(T)}{\Var(X)} \;=\; \frac{\Var(T)}{\Var(T)+\Var(E)},
\]
estimated in practice by the correlation (for continuous scores) or
percent agreement / $\kappa$ (for categorical labels) between two
independent applications of the same measurement procedure to the same
items --- exactly Benchmark~5's construction: Qwen labels the same $101$
clips three times with sampling ($T{=}0.7$), and the pooled pairwise
agreement across the $\binom{3}{2}=3$ run-pairs estimates $\rho_{XX'}$
for each field.

\begin{table}[h]
\centering
\begin{tabular}{lrr}
\toprule
Field & Qwen self-agreement $\hat\rho_{XX'}$ & human$\leftrightarrow$human agreement \\
\midrule
scenario  & 88.1\% & 64.8\% \\
verb      & 91.2\% & 69.5\% \\
condition & 90.4\% & 71.5\% \\
flavor    & 93.4\% & 64.6\% \\
noun      & 87.5\% & 60.7\% \\
\bottomrule
\end{tabular}
\caption{Benchmark 5, Part A: Qwen is more self-consistent than independent
humans are with each other, on every field.}
\end{table}

\subsection{The attenuation bound: reliability limits, but does not imply, validity}

Let $Y$ denote the (human-consensus) ground truth and $X$ the model's
label, both modeled as noisy measurements of a shared underlying construct.
The classical \emph{attenuation formula} relates the correlation between
two noisy measurements to their individual reliabilities:
\[
  \rho_{XY} \;\le\; \sqrt{\rho_{XX'}\,\rho_{YY'}}.
\]

\begin{proposition}[Reliability is necessary, not sufficient, for validity]
\label{prop:necessary-not-sufficient}
High self-agreement $\rho_{XX'}$ upper-bounds, but does not lower-bound,
agreement with ground truth $\rho_{XY}$: $\rho_{XX'}=1$ is consistent with
$\rho_{XY}=0$.
\end{proposition}
\begin{proof}
Let $X \equiv b$ be a \emph{constant} predictor (formally, $X_i = b$ for
every item $i$, independent of any input). Two independent applications of
this ``measurement procedure'' agree on every item, so its test--retest
agreement is $1$ (perfect reliability: $\Var(E) = 0$, so $\rho_{XX'}=1$
trivially, since $\Var(T) = \Var(X) = 0$ degenerates the ratio to $1$ by
convention for a constant). Yet if the true label $Y_i$ varies across items
and is not itself constant, $\Cov(X,Y) = 0$ identically (a constant has
zero covariance with anything), so $\rho_{XY} = 0$. This satisfies the
attenuation bound ($0 \le \sqrt{1\cdot\rho_{YY'}}$) with equality only
failing to be tight, exhibiting a case of maximal reliability and zero
validity.
\end{proof}

\begin{corollary}[Application to the observed self-fed prior]
\label{cor:self-fed}
Model the \texttt{verb} agency sub-field, restricted to clips where the
true agency is \texttt{human-fed}, as $X_i \approx \texttt{self}$ for
(nearly) every such clip $i$ --- an approximately constant predictor on
this subpopulation, per Benchmark~2's finding of $0\%$ agreement with
truth across five independent elicitation conditions
(Table~\ref{tab:agency-conditions}, reproduced from Section~5). By
Proposition~\ref{prop:necessary-not-sufficient}'s mechanism, this
near-constant behavior is exactly what produces high measured
\emph{self-agreement} on the field overall ($91.2\%$, Table above,
pooling both agency subpopulations) while contributing systematically
zero validity on the human-fed subpopulation specifically. Reliability
computed on the \emph{marginal} field distribution therefore masks a
validity failure that is only visible once the field is stratified by
true label --- the same diagnostic principle underlying
Proposition~\ref{prop:difficulty}'s stratification for consensus accuracy.
\end{corollary}

\subsection{Bias--variance view of the same phenomenon}

Equivalently, write the per-item error as
$\Delta_i = \mathbb{1}[X_i \ne Y_i]$ and decompose the population error
rate on a subgroup $s$ (e.g.\ true label $=$ human-fed) as
\[
  \E_s[\Delta] \;=\; \underbrace{\big(\Pr(X = \texttt{self} \mid s) - 0\big)}_{\text{systematic bias on } s}
  \;+\; \underbrace{\text{(noise-driven disagreement with } Y \text{ within } s\text{)}}_{\text{variance term}} .
\]
A rater with zero variance (perfectly reliable, in the test--retest sense)
but $100\%$ systematic bias toward one label on a subgroup contributes
$\E_s[\Delta] = 1$ on exactly that subgroup regardless of how reliable it
is elsewhere --- reliability controls only the variance term, never the
bias term, which is precisely why five independent interventions that only
perturb \emph{how the question is asked} (Section~\ref{sec:significance})
cannot repair a bias-dominated failure mode: none of them supply the
missing visual/attentional signal that would move $\Pr(X=\texttt{self}\mid
\text{human-fed})$ away from its default.

\section{Significance of the human-fed blind spot across independent interventions}
\label{sec:significance}

\subsection{A single-condition test}

Benchmark~2 evaluated agency accuracy on the $n=17$ clips with true label
\texttt{human-fed} under several independent elicitation conditions:

\begin{table}[h]
\centering
\begin{tabular}{lr}
\toprule
Condition & Agency accuracy, true human-fed ($n=17$) \\
\midrule
Composite verb question                & $0\%$ \\
Agency asked alone                      & $0\%$ \\
``Whose hand holds the food?'' (perception probe) & $6\%$ \\
Describe-people-first scaffold          & $18\%$ \\
``fed-by-another-person'' rewording     & $0\%$ \\
\bottomrule
\end{tabular}
\caption{Reproduced from Benchmark 2 / Benchmark 5.}
\label{tab:agency-conditions}
\end{table}

Fix a null hypothesis $H_0$: the model's per-trial probability of correctly
identifying agency on a true human-fed clip is some baseline rate $p_0$
representative of ``the model can perceive agency at least as well as it
does on the easier self-fed subgroup'' (empirically $\approx 90\%$ there),
or, more conservatively, $p_0 = 0.5$ (``no better than an agency coin
flip''). Under $H_0$, with $K$ the number of correct calls out of $n=17$
i.i.d.\ trials, $K \sim \mathrm{Binomial}(17, p_0)$, and the observed
$K=0$ has one-sided p-value
\[
  \Pr(K \le 0 \mid n=17, p_0) \;=\; (1-p_0)^{17}.
\]
\[
  p_0 = 0.5: \quad (0.5)^{17} \approx 7.6\times 10^{-6}, \qquad
  p_0 = 0.9: \quad (0.1)^{17} \approx 10^{-17}.
\]
Either way, observing zero correct calls out of 17 is not plausibly
explained by ordinary sampling noise around a reasonable baseline rate ---
it is evidence the true per-trial success probability on this subgroup is
near zero, not merely below its self-fed-subgroup counterpart.

\subsection{Combining across independent conditions: Fisher's method}

Table~\ref{tab:agency-conditions} contains five \emph{semi-independent}
probes (independent in the sense that each perturbs a different aspect of
elicitation --- question framing, isolated sub-question, pure perception,
reasoning scaffold, label vocabulary --- rather than being repeated trials
of the identical prompt). Let $p_i$ be the one-sided binomial p-value
computed as above (under $H_0{:}\,p_0=0.5$) for condition $i \in
\{1,\dots,5\}$, using the reported accuracy to back out an equivalent
integer success count $K_i = \mathrm{round}(0.17 \times \mathrm{acc}_i)$
out of $n=17$:\footnote{$K_i = 0$ for the three $0\%$ conditions, $K_i=1$
for the $6\%$ probe, $K_i=3$ for the $18\%$ scaffold.}

\begin{definition}[Fisher's combined probability test]
\[
  \chi^2_{\mathrm{combined}} \;=\; -2\sum_{i=1}^{5} \ln p_i \;\sim\; \chi^2_{2\times 5 = 10} \quad \text{under } H_0.
\]
\end{definition}

With $p_1=p_2=p_5 = \Pr(K\le 0\mid n{=}17,p_0{=}0.5) \approx 7.63\times10^{-6}$
(the three $0\%$ conditions),
$p_3 = \Pr(K\le 1\mid n{=}17,p_0{=}0.5) \approx 1.37\times10^{-4}$, and
$p_4 = \Pr(K\le 3\mid n{=}17,p_0{=}0.5) \approx 6.36\times10^{-3}$,
\[
  \chi^2_{\mathrm{combined}} \;\approx\; -2\big(3\ln(7.63\times10^{-6}) + \ln(1.37\times10^{-4}) + \ln(6.36\times10^{-3})\big) \;\approx\; 98.6,
\]
against a $\chi^2_{10}$ null distribution whose $99.9$th percentile is
about $29.6$ --- the combined statistic exceeds it by more than $3\times$,
giving a combined p-value (numerically, $\mathrm{sf}_{\chi^2_{10}}(98.6)
\approx 1.0\times10^{-16}$) many orders of magnitude below any conventional
significance threshold.

\begin{remark}[What this rules in and rules out]
This computation does not claim the five conditions are statistically
independent in the strict i.i.d.\ sense (they share the same 17 clips and
the same underlying model weights), so the nominal $\chi^2_{10}$
$p$-value should be read qualitatively rather than taken at face value.
What it does formalize is the qualitative argument used in
Section~\ref{sec:reliability}: a fixable \emph{prompting} deficiency
predicts that \emph{at least one} of five substantially different framings
(especially the pure-perception probe, which asks nothing about the
taxonomy at all) should recover non-trivial accuracy, whereas a
\emph{perceptual} deficiency in how the visual encoder or attention
mechanism attributes a hand to a person predicts exactly the observed
pattern: uniformly near-zero accuracy regardless of how the question is
phrased, with only the scaffold condition --- which adds an explicit
step of \emph{describing} what is seen before answering --- producing a
partial (and still far from adequate) recovery to $18\%$.
\end{remark}

\subsection{Interaction with the few-shot result}

Benchmark~5, Part B adds a sixth and seventh condition (4-shot and 8-shot
in-context demonstrations, the latter explicitly including two human-fed
examples) and observes $0\%$ agency accuracy on true human-fed clips at
\emph{both} shot counts. Extending the binomial argument, providing
correctly labeled human-fed exemplars in-context is the single
intervention most directly targeted at a demonstration-availability
explanation of the failure (``the model has never been shown what
human-fed looks like in this schema''); its failure to move $K$ away from
$0$ is the strongest single piece of evidence in the suite that the
deficit is not remediable by prompting or in-context learning at the
current model scale, and would require either fine-tuning on labeled
human-fed examples or an architectural/attention-level intervention (e.g.\
person-conditioned attention, as sketched by the scaffold condition's
partial success) to resolve.


\end{document}
